% Figures du juge de la construction pour le Tome III (genome/JUGEMENT-construction.md,
% genome/construction-juge-pour-le-memoire.md). Gardé par \ifdefined\GrapheJconsCodon. Données :
% \Racine/figures/tome3/donnees/juge-*.csv et juge-nombres.tex, écrits par memoire/outils/tome3_juge.py depuis
% mesures/genome/construction/juge/ (et, pour le mode exact sur R1, mesures/moteur/O/journaux/o3/). Styles « jcons ».
%
% Second flux (2026-10-27) : dictionnaire des couleurs du Tome III (palette-tome3.tex) et légendes G9 (\Titre,
% \Legende, \Source ; le chapitre écrit \LegendeFigure{<Nom>}). III-8.8 : axe tiré des données (G2 du juge des
% chiffres, aucun expert hors du cadre) ; III-10.1 : D21 du juge de design ; III-10.3 : D22 et M13.
%   III-8.8 JconsCodon   III-10.1 JconsExperts   III-10.3 JconsFenetres

\ifdefined\GrapheJconsCodon\else
\input{\Racine/figures/tome3/legende-tome3}
\providecommand{\CodonDonnees}{\Racine/figures/tome3/donnees}
\providecommand{\CodonNum}[1]{\num[group-digits=integer]{#1}}
\input{\CodonDonnees/juge-nombres}

\pgfplotsset{
  % Catégories de la grille des experts (D21) : 0 aucune fenêtre (blanc, cadre et croix d'encre), [10, 20) R1 seul
  % (CouleurRun plein), [20, 30) neuves seules (CouleurNeuves plein, un point blanc), [30, 40) les deux, du gris
  % clair au gris moyen selon log10(1 + appels) : le bruit de fond ne prend plus l'œil aux exceptions. Troisième
  % passe (DN21 du juge de design) : la rampe se resserre, black!5 à black!20 au lieu de black!8 à black!35.
  colormap={jcons}{color(0)=(white) color(9.99)=(white)
    color(10)=(CouleurRun) color(19.99)=(CouleurRun)
    color(20)=(CouleurNeuves) color(29.99)=(CouleurNeuves)
    color(30)=(black!5) color(40)=(black!20)},
}
\tikzset{jcons note/.style={t3 sans cesure, font=\sffamily\scriptsize, text=black!70, align=center}}

% =============================================================================================================
% III-8.8 (F-J6). Le codon par le décodeur du juge.
\newcommand{\GrapheJconsCodon}{%
\begin{tikzpicture}
\begin{axis}[t3 graphe, t3 points, t3 decimales x=1, name=g, width=0.6\linewidth, height=6.2cm,
  title={(a) \EF{share of $-0$ in each drawn expert}{part des $-0$ de chaque expert tiré}},
  xlabel={\EF{relative depth (layer / last layer)}{profondeur relative (couche / dernière couche)}},
  ylabel={\EF{$-0$ codes (\%)}{codes $-0$ (\%)}}, xmin=-0.03, xmax=1.03, ymin=\JconsNPartAxeMin, ymax=\JconsNPartAxeMax,
  legend style={at={(0.98,0.86)}, anchor=north east}, legend columns=1]
\draw[CouleurVingtB, dashed, line width=0.7pt] (axis cs:-0.03,\JconsNPartVingt) -- (axis cs:1.03,\JconsNPartVingt);
\draw[CouleurCentVingtB, dash pattern=on 1.5pt off 1.5pt, line width=0.7pt] (axis cs:-0.03,\JconsNPartCent) -- (axis cs:1.03,\JconsNPartCent);
\addplot[only marks, t3 vingtb point, fill opacity=0.75, unbounded coords=discard]
  table[x=x, y expr={\thisrow{taille}==20 ? \thisrow{part} : nan}] {\CodonDonnees/juge-codon-experts.csv};
\addlegendentry{gpt-oss-20b (\JconsNExpertsVingt)}
\addplot[only marks, t3 centvingtb point, fill opacity=0.75, unbounded coords=discard]
  table[x=x, y expr={\thisrow{taille}==120 ? \thisrow{part} : nan}] {\CodonDonnees/juge-codon-experts.csv};
\addlegendentry{gpt-oss-120b (\JconsNExpertsCent)}
\node[t3 nombre=CouleurCentVingtB!85!black, anchor=west] at (axis cs:0.02,\JconsNPartMax)
  {\JconsNPartMaxModele, \EF{layer}{couche} \JconsNPartMaxCouche, expert \JconsNPartMaxExpert\Dp\ \Pourcent{\JconsNPartMax}};
\addlegendimage{CouleurVingtB, dashed, line width=0.7pt}
\addlegendentry{\EF{whole 20b file}{fichier 20b entier}\Dp\ \Pourcent{\JconsNPartVingt}}
\addlegendimage{CouleurCentVingtB, dash pattern=on 1.5pt off 1.5pt, line width=0.7pt}
\addlegendentry{\EF{whole 120b file}{fichier 120b entier}\Dp\ \Pourcent{\JconsNPartCent}}
\node[t3 note, anchor=south west] at (axis cs:0.0,\JconsNPartAxeMin+0.15)
  {\EF{nibbles 8 in the 96 decodings:}{demi-octets 8 dans les 96 décodages :} \JconsNHuit};
\end{axis}
\begin{axis}[t3 graphe, t3 decimales y=2, name=d, at={(g.east)}, anchor=west, xshift=1.55cm, width=0.36\linewidth, height=6.2cm,
  title={(b) \EF{gain per expert weight}{gain par poids d'expert}},
  ylabel={\EF{gain (bit/weight)}{gain (bit/poids)}}, ymin=0, ymax=0.16, xmin=-0.6, xmax=1.6, xtick={0,1}, xmajorgrids=false,
  xticklabels={20b, 120b}, clip=false]
\node[jcons note, anchor=north, yshift=-0.42cm] at (axis cs:0,0)
  {\CodonNum{\JconsNGoVingtA}\\$\downarrow$\\\CodonNum{\JconsNGoVingtB}~\EF{GB}{Go}};
\node[jcons note, anchor=north, yshift=-0.42cm] at (axis cs:1,0)
  {\CodonNum{\JconsNGoCentA}\\$\downarrow$\\\CodonNum{\JconsNGoCentB}~\EF{GB}{Go}};
\draw[fill=CouleurVingtB, draw=CouleurVingtB!70!black, line width=0.4pt] (axis cs:-0.3,0) rectangle (axis cs:0.3,\JconsNGainVingt);
\draw[fill=CouleurCentVingtB, draw=CouleurCentVingtB!70!black, line width=0.4pt] (axis cs:0.7,0) rectangle (axis cs:1.3,\JconsNGainCent);
\draw[t3 seuil] (axis cs:-0.6,\JconsNSeuilCdeux) -- (axis cs:1.6,\JconsNSeuilCdeux);
\node[t3 etiquette, anchor=south] at (axis cs:0.5,0.142) {A-C2\Dp\ $\ge\CodonNum{\JconsNSeuilCdeux}$};
\node[t3 nombre=CouleurVingtB, anchor=north] at (axis cs:0,0.105) {\CodonNum{\JconsNGainVingt}};
\node[t3 nombre=CouleurCentVingtB!85!black, anchor=north] at (axis cs:1,0.105) {\CodonNum{\JconsNGainCent}};
\end{axis}
\end{tikzpicture}}
\Titre{JconsCodon}{The codon by the judge's decoder}{Le codon par le décodeur du juge}
\Legende{JconsCodon}{\mes}
  {The judge's own decoder finds the canonical expert in all \JconsNExpertsTient\ drawn experts, and the gain of
   both files passes the threshold declared before.}
  {Le décodeur propre du juge retrouve l'expert canonique dans les \JconsNExpertsTient\ experts tirés, et le gain
   des deux fichiers passe le seuil déclaré avant.}
  {(a) The \JconsNExperts\ experts drawn from the two v1.2 files and decoded by the judge: share of $-0$ codes in
   the GGUF expert, by relative depth; dashed, the share over the whole file. All \JconsNExpertsTient\ decode to
   the canonical expert, with no nibble 8 left; the \JconsNTemoins\ v1.0 controls decode to the GGUF bytes
   (\JconsNTemoinsTient/\JconsNTemoins). The set of nibbles changed between v1.0 and v1.2 is exactly the set of
   $-0$ (note of the judge §~1.2). The share of $-0$ varies by expert: the highest, \Pourcent{\JconsNPartMax},
   is in layer~\JconsNPartMaxCouche\ of the \JconsNPartMaxModele; the file gain does not depend on it. The axis is
   drawn from the data: no expert lies outside the frame. (b) The gain per expert weight of the four files read
   whole (aligned chunks), their sizes below, against the threshold A-C2 declared before.}
  {(a) Les \JconsNExperts\ experts tirés des deux fichiers v1.2 et décodés par le juge : part des codes $-0$ de
   l'expert du GGUF, selon la profondeur relative ; en tirets, la part sur le fichier entier. Les
   \JconsNExpertsTient\ se décodent en l'expert canonique, sans aucun demi-octet 8 ; les \JconsNTemoins\ témoins
   v1.0 se décodent en les octets du GGUF (\JconsNTemoinsTient/\JconsNTemoins). L'ensemble des demi-octets changés
   entre v1.0 et v1.2 est exactement celui des $-0$ (note du juge §~1.2). La part des $-0$ varie d'un expert à
   l'autre : la plus haute, \Pourcent{\JconsNPartMax}, est à la couche~\JconsNPartMaxCouche\ du
   \JconsNPartMaxModele ; le gain du fichier n'en dépend pas. L'axe est tiré des données : aucun expert n'est hors
   du cadre. (b) Le gain par poids d'expert des quatre fichiers lus en entier (tronçons alignés), leurs tailles
   en dessous, contre le seuil A-C2 déclaré avant.}
\Source{JconsCodon}{\fichier{mesures/genome/construction/juge/J-T1-echantillon.csv}, \fichier{J-T0-fichiers.json}\Pv{}
  \fichier{mesures/genome/construction/codon/tailles.json}\Pv\ \EF{script}{script} \fichier{memoire/outils/tome3_juge.py}.}

% =============================================================================================================
% III-10.1 (F-J4). Les experts vus : grille 24 × 32 du 20b, qui les a routés et combien.
% Cases carrées : axe de 12,6 cm sur 9,45 cm, 32 × 24 cases de 0,394 cm.
\newcommand{\JconsCle}[3]{% x, remplissage, texte : une entrée de la clé
  \fill[#2] (#1,0) rectangle ++(0.36,0.36);\draw[black!40, line width=0.3pt] (#1,0) rectangle ++(0.36,0.36);
  \node[jcons note, anchor=west] at (#1+0.42,0.18) {#3};}
\newcommand{\GrapheJconsExperts}{%
\begin{tikzpicture}
\begin{axis}[name=grille, scale only axis, width=12.6cm, height=9.45cm, enlargelimits=false, clip=false,
  xmin=-0.5, xmax=31.5, ymin=-0.5, ymax=23.5, y dir=reverse,
  xtick={0,4,...,28,31}, ytick={0,2,...,22}, tick label style={font=\scriptsize}, tick style={draw=none},
  xlabel={\EF{expert}{expert}}, ylabel={\EF{layer}{couche}}, label style={font=\small},
  colormap name=jcons, point meta min=0, point meta max=40, axis line style={black!40}]
\addplot[matrix plot*, mesh/cols=32, point meta=explicit] table[x=expert, y=couche, meta=m] {\CodonDonnees/juge-grille.csv};
\addplot[only marks, mark=*, mark size=1.1pt, mark options={solid, fill=white, draw=white}, unbounded coords=discard]
  table[x=expert, y expr={\thisrow{sorte}==2 ? \thisrow{couche} : nan}] {\CodonDonnees/juge-grille.csv};
\addplot[only marks, mark=square, mark size=5.2pt, mark options={solid, draw=Encre, line width=0.6pt}, unbounded coords=discard]
  table[x=expert, y expr={\thisrow{sorte}==0 ? \thisrow{couche} : nan}] {\CodonDonnees/juge-grille.csv};
\addplot[only marks, t3 cassure point, mark size=3.2pt, unbounded coords=discard]
  table[x=expert, y expr={\thisrow{sorte}==0 ? \thisrow{couche} : nan}] {\CodonDonnees/juge-grille.csv};
% le compte des experts jamais routés, couche par couche, en marge
\addplot[only marks, mark=none, nodes near coords, point meta=explicit symbolic,
  nodes near coords style={font=\sffamily\scriptsize, anchor=center, text=Encre}]
  table[x expr=32.4, y=couche, meta=vides] {\CodonDonnees/juge-grille-marge.csv};
\node[jcons note, anchor=south, text=Encre] at (axis cs:32.4,-0.6) {$\times$};
\end{axis}
\begin{scope}[shift={($(grille.south west)+(0,-1.25)$)}]
\JconsCle{0}{CouleurRun}{\EF{R1 only}{R1 seul} (\JconsNRunSeul)}
\JconsCle{2.7}{CouleurNeuves}{\EF{new windows only}{fenêtres neuves seules} (\JconsNNeufSeul)}
\fill[white] (2.88,0.18) circle (1.1pt);
\shade[left color=black!5, right color=black!20] (7.2,0) rectangle ++(0.9,0.36);\draw[black!40, line width=0.3pt] (7.2,0) rectangle ++(0.9,0.36);
\node[jcons note, anchor=west] at (8.16,0.18) {\EF{both}{les deux} (\JconsNDeux)};
\draw[Encre, line width=0.6pt] (10.1,0) rectangle ++(0.36,0.36);
\path[t3 cassure point, mark size=2.4pt] plot coordinates {(10.28,0.18)};
\node[jcons note, anchor=west] at (10.52,0.18) {\EF{never routed}{jamais routé} (\JconsNVides)};
\end{scope}
\end{tikzpicture}}
\Titre{JconsExperts}{The experts that were seen}{Les experts vus}
\Legende{JconsExperts}{\mes}
  {\JconsNVides\ experts of gpt-oss-20b were never routed, on either set of windows: no KL measurement of the
   construction called them (their decoding, for its part, is checked by value: C1, \JconsNExpertsDecodes\ experts
   of \JconsNExpertsTous).}
  {\JconsNVides\ experts de gpt-oss-20b n'ont jamais été routés, sur aucun des deux jeux de fenêtres : aucune
   mesure de KL de la construction ne les a appelés (leur décodage, lui, est vérifié en valeur : C1,
   \JconsNExpertsDecodes\ experts sur \JconsNExpertsTous).}
  {The $24\times32$ experts of gpt-oss-20b, each cell coloured by the windows that routed tokens to it: the ten
   windows of R1 only (blue), the ten new windows of the judge only (purple, white dot), both (grey, darker with
   the number of calls on a logarithmic scale), or none (framed and crossed; their count per layer in the margin).
   A call is one token routed to an expert; a window routes its \JconsNJetonsFenetre\ tokens, of which
   \JconsNContexte\ are context not scored (\CodonNum{\JconsNAppelsRun} calls in all on each set of windows; at most
   \CodonNum{\JconsNAppelsMaxUnJeu} for one expert on one set, \CodonNum{\JconsNAppelsMax} on both together). It is
   the map of the space the router visits on real text, under the routing of the reference: a change confined to
   the \JconsNVides\ crossed experts would move no KL of \crefrange{t3cod:ch}{t3jug:ch} measured with that routing, but a variant
   whose routing flips could call them.}
  {Les $24\times32$ experts de gpt-oss-20b, chaque case colorée selon les fenêtres qui lui ont routé des jetons :
   les dix fenêtres de R1 seules (bleu), les dix fenêtres neuves du juge seules (violet, point blanc), les deux
   (gris, plus foncé avec le nombre d'appels, en échelle logarithmique), ou aucune (encadrée et barrée ; leur compte
   par couche en marge). Un appel est un jeton routé vers un expert ; une fenêtre route ses \JconsNJetonsFenetre\
   jetons, dont \JconsNContexte\ de contexte non notés (\CodonNum{\JconsNAppelsRun} appels en tout sur chaque jeu
   de fenêtres ; au plus \CodonNum{\JconsNAppelsMaxUnJeu} pour un expert sur un jeu, \CodonNum{\JconsNAppelsMax} sur
   les deux réunis). C'est la carte de l'espace que le routeur parcourt sur du texte réel, sous le routage de la
   référence : un changement confiné aux \JconsNVides\ experts barrés ne bougerait aucune KL des \crefrange{t3cod:ch}{t3jug:ch}
   mesurée avec ce routage, mais une variante dont le routage bascule pourrait les appeler.}
\Source{JconsExperts}{\fichier{mesures/genome/construction/juge/figures/routage-grille.csv},
  \fichier{J-T0-fichiers.json} ($k$)\Pv\ \fichier{mesures/genome/construction/synonyme/kl/kl-frontiere-768-positions.csv}
  (\EF{scored positions}{positions notées})\Pv\ \fichier{mesures/genome/construction/codon/tailles.json}
  (C1)\Pv\ \EF{script}{script} \fichier{memoire/outils/tome3_juge.py}.}

% =============================================================================================================
% III-10.3 (F-J7). La difficulté des fenêtres : KL par fenêtre, R1 et neuves côte à côte.
\newcommand{\JconsSerie}[4]{% colonne, style du trait, style de la marque, entrée de légende
  \addplot[#2, line width=0.5pt, #3, unbounded coords=jump] table[x=x, y=#1] {\CodonDonnees/juge-kl-fenetres.csv};
  \addlegendentry{#4}}
\newcommand{\GrapheJconsFenetres}{%
\begin{tikzpicture}
\begin{axis}[t3 graphe, width=0.95\linewidth, height=7cm, ymode=log, ymin=5e-5, ymax=1,
  xmin=-0.6, xmax=20.6, xtick={0,...,9,11,12,...,20}, xticklabels={0,...,9,0,1,...,9},
  xticklabel style={font=\scriptsize}, xlabel={\EF{window}{fenêtre}}, xmajorgrids=false,
  ylabel={\EF{mean KL of the window}{KL moyenne de la fenêtre}},
  legend style={at={(0.5,1.03)}, anchor=south}, legend columns=4]
\fill[CouleurNeuves!8] (axis cs:10.5,5e-5) rectangle (axis cs:20.6,1);
\draw[black!40, line width=0.5pt] (axis cs:10,5e-5) -- (axis cs:10,1);
\JconsSerie{frontiere}{CouleurFrontiere, solid}{t3 frontiere point}{\EF{frontier}{frontière}}
\JconsSerie{lambda1}{CouleurSynonyme, dash pattern=on 3pt off 1.5pt}{t3 synonyme point}{$\lambda=1$}
\JconsSerie{lambda0}{CouleurLambdaZero, dash pattern=on 3pt off 1pt on 1pt off 1pt}{t3 lambda zero point}{$\lambda=0$}
\JconsSerie{exact}{CouleurExact, solid}{t3 exact point}{\EF{exact mode}{mode exact}}
\node[t3 etiquette, anchor=south] at (axis cs:4.5,6e-5) {R1};
\node[t3 etiquette, anchor=south] at (axis cs:15.5,6e-5) {\EF{new windows (judge)}{fenêtres neuves (juge)}};
\end{axis}
\end{tikzpicture}}
\Titre{JconsFenetres}{The difficulty of the windows}{La difficulté des fenêtres}
\Legende{JconsFenetres}{\mes}
  {Some windows are harder than others for every variant alike; the exact mode stays far below them all.}
  {Certaines fenêtres sont plus difficiles que d'autres pour toutes les variantes à la fois ; le mode exact reste
   loin au-dessous de toutes.}
  {Mean KL of each window (\JconsNNotees\ scored positions) for the frontier, the two synonyms and the exact mode,
   on the ten windows of R1 and the ten new windows of the judge (shaded). The thin lines join windows that do not
   follow one another: they only show that the variants rise and fall together. Paths: on R1, the variants
   through the substituted f64 reference and the exact mode from the per-position file of lot O; on the new
   windows, $\lambda=0$, $\lambda=1$ and the exact mode through the engine, the frontier through the substituted
   reference. One path is drawn per series and per set. On R1 the three variants also went through the engine:
   engine minus substituted reference, \CodonNum{\JconsNCheminRunMoy} on average for the frontier,
   \CodonNum{+\JconsNCheminRunMoyUn} for $\lambda=1$, \CodonNum{\JconsNCheminRunMoyZero} for $\lambda=0$, from
   \CodonNum{\JconsNCheminRunMin} to \CodonNum{+\JconsNCheminRunMax} by window over the three; on the new windows,
   $\lambda=0$ went through both, \CodonNum{\JconsNCheminNeufMoy} on average. The
   frontier ranges from \CodonNum{\JconsNFrontMin} to \CodonNum{\JconsNFrontMax} on R1: this spread from window to
   window, shared by all variants, makes the standard deviation by windows ($\sigma_f=\CodonNum{\JconsNSigmaF}$ for
   the frontier) much larger than the naive one by positions, and pairing removes it (\cref{fig:t3-syno-foret}).
   The new windows are drawn from wiki.test, the corpus of R1.}
  {KL moyenne de chaque fenêtre (\JconsNNotees\ positions notées) pour la frontière, les deux synonymes et le mode
   exact, sur les dix fenêtres de R1 et les dix fenêtres neuves du juge (fond teinté). Les traits fins joignent des
   fenêtres qui ne se suivent pas : ils ne montrent que le parallélisme des variantes. Chemins : sur R1, les
   variantes par la référence f64 substituée et le mode exact depuis le fichier par position du lot O ; sur les
   fenêtres neuves, $\lambda=0$, $\lambda=1$ et le mode exact par le moteur, la frontière par la référence
   substituée. Un seul chemin est tracé par série et par jeu. Sur R1, les trois variantes sont aussi passées par
   le moteur : moteur moins référence substituée, \CodonNum{\JconsNCheminRunMoy} en moyenne pour la frontière,
   \CodonNum{+\JconsNCheminRunMoyUn} pour $\lambda=1$, \CodonNum{\JconsNCheminRunMoyZero} pour $\lambda=0$, de
   \CodonNum{\JconsNCheminRunMin} à \CodonNum{+\JconsNCheminRunMax} par fenêtre sur les trois ; sur les fenêtres
   neuves, $\lambda=0$ est passé par les deux, \CodonNum{\JconsNCheminNeufMoy} en moyenne. La frontière va de \CodonNum{\JconsNFrontMin} à \CodonNum{\JconsNFrontMax} sur R1 : cet écart
   d'une fenêtre à l'autre, commun à toutes les variantes, rend l'écart-type par fenêtres
   ($\sigma_f=\CodonNum{\JconsNSigmaF}$ pour la frontière) bien plus grand que l'écart-type naïf par positions, et
   l'appariement le retire (\cref{fig:t3-syno-foret}). Les fenêtres neuves sont tirées de wiki.test, le corpus de R1.}
\Source{JconsFenetres}{\fichier{mesures/genome/construction/juge/figures/kl-par-fenetre.csv}\Pv{}
  \fichier{mesures/moteur/O/journaux/o3/c4-20b-comparer.csv}\Pv{}
  \fichier{mesures/genome/construction/synonyme/moteur/kl-moteur-frontiere-positions.csv}\Pv\ \EF{script}{script}
  \fichier{memoire/outils/tome3_juge.py}.}
\fi
